%! Author = chtompki
%! Date = 1/14/26
\documentclass[12pt]{amsart}
% Default margins are too wide all the way around. I reset them here
\setlength{\hoffset}{-.75in}
\setlength{\textwidth}{6.5in}
\setlength{\voffset}{-.5in}
\setlength{\textheight}{9.0in}
\usepackage{amsmath}
\usepackage{amssymb}
\usepackage{cancel}
\usepackage{amsthm}
\usepackage{enumitem}
\usepackage{setspace}
\usepackage{graphicx}
\usepackage{hyperref}

\renewcommand{\thesubsection}{\thesection.\alph{subsection}}

\newenvironment{solution}
{\begin{proof}[Solution]}
{\end{proof}}

\NewDocumentCommand{\vitem}{v}{%
    \item[\texttt{#1}]%
}

\title{Numerical Analysis\\Problem Set 5}
\author{Rob Tompkins}
\date{\today}

\begin{document}
    \maketitle
    \date{\today}
    \begin{onehalfspace}

        % Problem 1
        \section{Problem 1}
        Derive the trapezoid integration scheme as an Adams-Moulton method with $m = 1$ step:
        \[
            \frac{w_{j+1}-w_j}{h} = \frac{1}{2}[f(t_{j+1},w_{j+1}) + f(t_j, w_j)]
        \]
        to approximate the initial value problem $y'=f(t,y)$ for $a\le t \le b$ with initial condition
        $y(a)=\alpha$.
        Also derive the associated error term.
        Let $w_j \approx y_j := y(t_j)$.
        \begin{itemize}
            \item[(a)] Use the Lagrange form of the polynomial interpolant to write
            \[
                f(t, y(t)) = P_1(t) + R_1(t).
            \]
            Integrate $P_1(t) = L_{1,0}(t)f(t_{j+1}, y_{j+1}) + L_{1,1}(t)f(t_j,y_j)$ from $t = t_j$ to $t=t_j+1$.
            \item[(b)] Set $w_{j+1} = w_j + \int_{t_j}^{t_{j+1}} f(t,y(t))dt$, approximating the 2nd term with
            $\int_{t_j}^{t_{j+1}} P_1(t)dt$, and derive the method above.
            \item[(c)] The error term $R_1(t) = \frac{f''(\xi,y(\xi))}{2}(t-t_{j+1})(t-t_j)$ for $f$ (some
            $\xi\in(t_j,t_{j+1})$) can be integrated from $t=t_j$ to $t=t_{j+1}$ to get the error associated
            with the Adams-Moulton scheme:
            \[
                \text{Error} = Mh^3 f''(\mu, y(\mu)).
            \]
            Do this integration to find value of $M$.
            \item[(d)] What is the order of this one-step implicit method?
        \end{itemize}

        \begin{solution}
        \end{solution}

    \end{onehalfspace}
\end{document}